\documentclass[12pt,a4paper]{article}
\usepackage[T1]{fontenc}
\usepackage[utf8]{inputenc}
\usepackage{lmodern}
\usepackage{amsmath,amssymb,amsthm,mathtools}
\usepackage{geometry}
\geometry{margin=1in}
\usepackage{microtype}
\usepackage{hyperref}
\usepackage{xcolor}
\usepackage{enumitem}
\usepackage{fancyhdr}
\usepackage{bookmark}

\usepackage{geometry}
\geometry{margin=1in}
\usepackage{microtype}
\usepackage{enumitem}
\usepackage{hyperref}
\usepackage{array}
\usepackage{booktabs}
\usepackage{xcolor}

\usepackage{booktabs}
\usepackage{longtable}
\usepackage{enumitem}
\usepackage{hyperref}

\hypersetup{
	colorlinks=true,
	linkcolor=blue,
	urlcolor=blue,
	citecolor=blue
}

\hypersetup{
	colorlinks=true,
	linkcolor=blue!50!black,
	urlcolor=blue!50!black,
	citecolor=blue!50!black,
	pdftitle={Theory of Commutative Algebra 14},
	pdfauthor={},
	pdfsubject={Solved problems and theory notes in commutative algebra},
	pdfcreator={OpenAI}
}

\setlist[itemize]{topsep=4pt,itemsep=2pt,parsep=0pt,partopsep=0pt}
\setlength{\parskip}{0.55em}
\setlength{\parindent}{0pt}

\setcounter{secnumdepth}{2}


\setcounter{tocdepth}{2}

\pagestyle{fancy}


\fancyhf{}
\lhead{Theory of Commutative Algebra 14}
\rhead{\thepage}
\renewcommand{\headrulewidth}{0.4pt}

\newtheoremstyle{notespace}{6pt}{6pt}{\normalfont}{} {\bfseries}{.}{0.5em}{}
\theoremstyle{notespace}
\newtheorem{definition}{Definition}[section]
\newtheorem{remark}{Remark}[section]


\newcommand{\Spec}{\operatorname{Spec}}
\newcommand{\Frac}{\operatorname{Frac}}
\newcommand{\OO}{\mathcal O}
\newcommand{\kk}{\mathbf{k}}
\newcommand{\A}{\mathbb A}
\newcommand{\Z}{\mathbb Z}
\newcommand{\Q}{\mathbb Q}
\newcommand{\R}{\mathbb R}
\newcommand{\C}{\mathbb C}
\newcommand{\F}{\mathbb F}
\newcommand{\mf}{\mathfrak}
\newcommand{\m}{\mathfrak m}
\newcommand{\p}{\mathfrak p}

\theoremstyle{definition}

\newtheorem{example}[definition]{Example}



\DeclareMathOperator{\Ass}{Ass}
\DeclareMathOperator{\Ann}{Ann}
\DeclareMathOperator{\Supp}{Supp}
\DeclareMathOperator{\Hom}{Hom}
\DeclareMathOperator{\Min}{Min}
\DeclareMathOperator{\Ker}{Ker}


\DeclareMathOperator{\depth}{depth}
\DeclareMathOperator{\Tor}{Tor}
\DeclareMathOperator{\Ext}{Ext}

\DeclareMathOperator{\coker}{coker}
\DeclareMathOperator{\im}{im}
\newcommand{\Coker}{\operatorname{Coker}}



\newtheorem{proposition}[definition]{Proposition}
\newtheorem{theorem}[definition]{Theorem}
\newtheorem{corollary}[definition]{Corollary}


\newcommand{\kerf}{\operatorname{Ker}}

\newcommand{\pd}{\operatorname{pd}}
\newcommand{\rank}{\operatorname{rank}}
\newcommand{\syz}{\operatorname{Syz}}
\newcommand{\ZZ}{\mathbb{Z}}







\hypersetup{
	colorlinks=true,
	linkcolor=blue!60!black,
	urlcolor=blue!60!black,
	citecolor=blue!60!black
}

\title{\vspace{-1.5em}\bfseries Theory of Commutative Algebra 14}
\author{}
\date{\today}

\begin{document}
	\maketitle
	\thispagestyle{plain}
		\begin{center}
		\begin{minipage}{0.92\textwidth}
			Lecture notes with worked examples in commutative algebra
		\end{minipage}
	\end{center}
		\pdfbookmark[1]{Contents}{contents}\tableofcontents
	\clearpage

\section{Localization at a prime, exact sequences, and the module-theoretic properties of \(A_{\mathfrak p}\)}

Let \(A\) be a commutative ring and let \(\mathfrak p \in \Spec(A)\) be a prime ideal. Put
\[
S=A\setminus \mathfrak p.
\]
Then the localization of \(A\) at \(\mathfrak p\) is
\[
A_{\mathfrak p}=S^{-1}A.
\]
Its elements are fractions
\[
\frac{a}{s},
\qquad a\in A,\ s\notin \mathfrak p.
\]

\subsubsection{Definition of the natural map \(A\to A_{\mathfrak p}\)}

There is a canonical ring homomorphism
\[
\iota:A\longrightarrow A_{\mathfrak p},
\qquad
a\longmapsto \frac{a}{1}.
\]
Its kernel is
\[
\ker(\iota)
=
\{a\in A:\exists s\notin \mathfrak p \text{ such that } sa=0\}.
\]
Hence one always has an exact sequence
\[
0
\longrightarrow
\ker(\iota)
\longrightarrow
A
\overset{\iota}{\longrightarrow}
A_{\mathfrak p}
\longrightarrow
\operatorname{coker}(\iota)
\longrightarrow
0.
\]

If \(A\) is a domain, then \(\iota\) is injective, so we get
\[
0
\longrightarrow
A
\longrightarrow
A_{\mathfrak p}
\longrightarrow
A_{\mathfrak p}/A
\longrightarrow
0.
\]

\subsubsection{Localization of a short exact sequence}

If \(I\subseteq A\) is an ideal, then there is the standard short exact sequence
\[
0
\longrightarrow
I
\longrightarrow
A
\longrightarrow
A/I
\longrightarrow
0.
\]
Localization preserves exactness, so after localizing at \(\mathfrak p\) we obtain
\[
0
\longrightarrow
I_{\mathfrak p}
\longrightarrow
A_{\mathfrak p}
\longrightarrow
(A/I)_{\mathfrak p}
\longrightarrow
0,
\]
where
\[
I_{\mathfrak p}=S^{-1}I.
\]

This is one of the most important exact sequences involving \(A_{\mathfrak p}\).

\subsubsection{The maximal ideal and the residue field}

The ring \(A_{\mathfrak p}\) is local. Its unique maximal ideal is
\[
\mathfrak p A_{\mathfrak p}.
\]
Therefore there is a natural short exact sequence
\[
0
\longrightarrow
\mathfrak p A_{\mathfrak p}
\longrightarrow
A_{\mathfrak p}
\longrightarrow
A_{\mathfrak p}/\mathfrak p A_{\mathfrak p}
\longrightarrow
0.
\]
The quotient
\[
\kappa(\mathfrak p):=A_{\mathfrak p}/\mathfrak p A_{\mathfrak p}
\]
is called the \emph{residue field} at \(\mathfrak p\). Thus the sequence may be written as
\[
0
\longrightarrow
\mathfrak p A_{\mathfrak p}
\longrightarrow
A_{\mathfrak p}
\longrightarrow
\kappa(\mathfrak p)
\longrightarrow
0.
\]

\subsubsection{Examples}

\paragraph{Example 1: \(A=\mathbb Z\), \(\mathfrak p=(5)\).}

We have
\[
\mathbb Z_{(5)}
=
\left\{
\frac{a}{b}\in \mathbb Q : 5\nmid b
\right\}.
\]
The local ring sequence is
\[
0
\longrightarrow
5\mathbb Z_{(5)}
\longrightarrow
\mathbb Z_{(5)}
\longrightarrow
\mathbb F_{5}
\longrightarrow
0.
\]

Also the canonical map
\[
\mathbb Z\longrightarrow \mathbb Z_{(5)}
\]
is injective, so there is an exact sequence
\[
0
\longrightarrow
\mathbb Z
\longrightarrow
\mathbb Z_{(5)}
\longrightarrow
\mathbb Z_{(5)}/\mathbb Z
\longrightarrow
0.
\]

\paragraph{Example 2: \(A=k[x]\), \(\mathfrak p=(x)\).}

Then
\[
A_{\mathfrak p}=k[x]_{(x)}.
\]
Taking \(I=(x)\), the localized quotient sequence is
\[
0
\longrightarrow
(x)k[x]_{(x)}
\longrightarrow
k[x]_{(x)}
\longrightarrow
(k[x]/(x))_{(x)}
\longrightarrow
0.
\]
Since
\[
k[x]/(x)\cong k,
\]
this becomes
\[
0
\longrightarrow
(x)k[x]_{(x)}
\longrightarrow
k[x]_{(x)}
\longrightarrow
k
\longrightarrow
0.
\]

\paragraph{Example 3: \(A=k[x,y]\), \(\mathfrak p=(x,y)\).}

Then
\[
A_{\mathfrak p}=k[x,y]_{(x,y)}.
\]
Its maximal ideal is
\[
(x,y)A_{\mathfrak p},
\]
and the residue field is
\[
A_{\mathfrak p}/(x,y)A_{\mathfrak p}\cong k.
\]
Hence
\[
0
\longrightarrow
(x,y)A_{\mathfrak p}
\longrightarrow
A_{\mathfrak p}
\longrightarrow
k
\longrightarrow
0.
\]

\paragraph{Example 4: \(A=\mathbb Z/6\mathbb Z\), \(\mathfrak p=(2)\).}

Here localization is more subtle because zero divisors are present. In this case some nonzero elements may become zero in \(A_{\mathfrak p}\). Thus the map
\[
A\longrightarrow A_{\mathfrak p}
\]
need not be injective. This gives a concrete example where the kernel of
\[
A\to A_{\mathfrak p}
\]
is nonzero.

\subsubsection{Abelian-group structure}

As with every \(A\)-module, the localization \(A_{\mathfrak p}\) is an abelian group under addition. Explicitly,
\[
\frac{a}{s}+\frac{b}{t}=\frac{at+bs}{st}.
\]
Thus, viewed only additively, \(A_{\mathfrak p}\) is always an abelian group.

\subsubsection{Is \(A_{\mathfrak p}\) free as an \(A\)-module?}

Usually, no. In general \(A_{\mathfrak p}\) is \emph{not} a free \(A\)-module.

For example, \(\mathbb Z_{(5)}\) is not a free \(\mathbb Z\)-module. Indeed, it is not cyclic as a \(\mathbb Z\)-module, because \(1\) does not generate \(\frac12\), and it does not admit a basis over \(\mathbb Z\).

So one should remember:

\[
A_{\mathfrak p}\ \text{is generally not free over }A.
\]

\subsubsection{Flatness of \(A_{\mathfrak p}\)}

Localization is the basic example of a flat module. In fact,
\[
A_{\mathfrak p}\ \text{is flat as an }A\text{-module.}
\]
Equivalently, for every short exact sequence of \(A\)-modules
\[
0
\longrightarrow
M'
\longrightarrow
M
\longrightarrow
M''
\longrightarrow
0,
\]
tensoring with \(A_{\mathfrak p}\) preserves exactness:
\[
0
\longrightarrow
M'\otimes_A A_{\mathfrak p}
\longrightarrow
M\otimes_A A_{\mathfrak p}
\longrightarrow
M''\otimes_A A_{\mathfrak p}
\longrightarrow
0.
\]

Moreover,
\[
M\otimes_A A_{\mathfrak p}\cong M_{\mathfrak p}
\]
for every \(A\)-module \(M\).

This explains why the sequence
\[
0
\longrightarrow
I
\longrightarrow
A
\longrightarrow
A/I
\longrightarrow
0
\]
localizes to
\[
0
\longrightarrow
I_{\mathfrak p}
\longrightarrow
A_{\mathfrak p}
\longrightarrow
(A/I)_{\mathfrak p}
\longrightarrow
0.
\]

\subsubsection{Faithfulness of \(A_{\mathfrak p}\)}

An \(A\)-module \(M\) is called \emph{faithful} if
\[
\operatorname{Ann}_A(M)=0.
\]
For \(A_{\mathfrak p}\), one computes
\[
\operatorname{Ann}_A(A_{\mathfrak p})
=
\{a\in A:\exists s\notin \mathfrak p \text{ such that } sa=0\}.
\]
Therefore \(A_{\mathfrak p}\) is faithful if and only if the canonical map
\[
A\longrightarrow A_{\mathfrak p}
\]
is injective.

In particular, if \(A\) is a domain, then \(A_{\mathfrak p}\) is faithful.

\paragraph{Faithful examples.}
\[
\mathbb Z_{(5)} \text{ is faithful as a } \mathbb Z\text{-module,}
\]
because
\[
\mathbb Z\longrightarrow \mathbb Z_{(5)}
\]
is injective.

Also,
\[
k[x]_{(x)} \text{ is faithful as a } k[x]\text{-module.}
\]

\paragraph{Non-faithful example.}
If \(A=\mathbb Z/6\mathbb Z\) and \(\mathfrak p=(2)\), then localization kills some nonzero elements, so \(A_{\mathfrak p}\) is not faithful as an \(A\)-module.

\subsubsection{Summary}

For a prime ideal \(\mathfrak p\subseteq A\), the localization
\[
A_{\mathfrak p}=(A\setminus \mathfrak p)^{-1}A
\]
fits naturally into several important exact sequences:
\[
0
\longrightarrow
I_{\mathfrak p}
\longrightarrow
A_{\mathfrak p}
\longrightarrow
(A/I)_{\mathfrak p}
\longrightarrow
0,
\]
\[
0
\longrightarrow
\mathfrak p A_{\mathfrak p}
\longrightarrow
A_{\mathfrak p}
\longrightarrow
\kappa(\mathfrak p)
\longrightarrow
0,
\]
and, when \(A\to A_{\mathfrak p}\) is injective,
\[
0
\longrightarrow
A
\longrightarrow
A_{\mathfrak p}
\longrightarrow
A_{\mathfrak p}/A
\longrightarrow
0.
\]

As an \(A\)-module, \(A_{\mathfrak p}\) is:

\[
\text{abelian: always,}
\]
\[
\text{free: usually not,}
\]
\[
\text{flat: always,}
\]
\[
\text{faithful: exactly when } A\to A_{\mathfrak p} \text{ is injective.}
\]

\subsection{Residue fields \(\kappa(\mathfrak p)\), stalks, and geometric meaning on \(\Spec(A)\)}

Let \(A\) be a commutative ring and let
\[
\mathfrak p\in \Spec(A)
\]
be a prime ideal. Associated to \(\mathfrak p\) are two fundamental local objects:

\[
A_{\mathfrak p},
\qquad
\kappa(\mathfrak p):=A_{\mathfrak p}/\mathfrak p A_{\mathfrak p}.
\]

The ring \(A_{\mathfrak p}\) is the local ring of \(A\) at the point \(\mathfrak p\), and the field \(\kappa(\mathfrak p)\) is called the \emph{residue field} at \(\mathfrak p\).

\subsubsection{1. Definition of the stalk \(A_{\mathfrak p}\)}

Recall that
\[
A_{\mathfrak p}=(A\setminus \mathfrak p)^{-1}A.
\]
Its elements are fractions
\[
\frac{a}{s},
\qquad
a\in A,\ s\notin \mathfrak p.
\]
The ring \(A_{\mathfrak p}\) is local, with unique maximal ideal
\[
\mathfrak p A_{\mathfrak p}.
\]

In algebraic geometry, if \(X=\Spec(A)\), then the structure sheaf \(\mathcal O_X\) satisfies
\[
\mathcal O_{X,\mathfrak p}\cong A_{\mathfrak p}.
\]
Thus \(A_{\mathfrak p}\) is the \emph{stalk} of the structure sheaf at the point \(\mathfrak p\).

\subsubsection{2. Definition of the residue field \(\kappa(\mathfrak p)\)}

Since \(A_{\mathfrak p}\) is a local ring with maximal ideal \(\mathfrak p A_{\mathfrak p}\), the quotient
\[
\kappa(\mathfrak p):=A_{\mathfrak p}/\mathfrak p A_{\mathfrak p}
\]
is a field. This is the residue field at \(\mathfrak p\).

There is a natural short exact sequence
\[
0
\longrightarrow
\mathfrak p A_{\mathfrak p}
\longrightarrow
A_{\mathfrak p}
\longrightarrow
\kappa(\mathfrak p)
\longrightarrow
0.
\]

Intuitively, \(\kappa(\mathfrak p)\) is obtained by taking all functions defined near the point \(\mathfrak p\), and then identifying two such functions if they differ by something vanishing at \(\mathfrak p\).

\subsubsection{3. Residue field as a field of fractions}

There is a very useful identification
\[
\kappa(\mathfrak p)\cong \operatorname{Frac}(A/\mathfrak p).
\]

Indeed, \(A/\mathfrak p\) is an integral domain because \(\mathfrak p\) is prime. Localizing \(A/\mathfrak p\) at its nonzero elements gives its field of fractions, and one checks that this is exactly
\[
A_{\mathfrak p}/\mathfrak p A_{\mathfrak p}.
\]

So one may think of \(\kappa(\mathfrak p)\) as the smallest field naturally attached to the point \(\mathfrak p\).

\subsubsection{4. Why \(A_{\mathfrak p}\) is called a stalk}

Let \(X=\Spec(A)\). For each basic open set
\[
D(f)=\{\mathfrak q\in \Spec(A):f\notin \mathfrak q\},
\]
the structure sheaf assigns the ring
\[
\mathcal O_X(D(f))=A_f.
\]

The stalk at \(\mathfrak p\) is the direct limit of all rings of sections over neighborhoods of \(\mathfrak p\):
\[
\mathcal O_{X,\mathfrak p}
=
\varinjlim_{\mathfrak p\in U}\mathcal O_X(U).
\]
In the affine case this direct limit is canonically isomorphic to
\[
A_{\mathfrak p}.
\]

So \(A_{\mathfrak p}\) contains the data of all regular functions defined in some neighborhood of the point \(\mathfrak p\), where two functions are identified if they agree on a smaller neighborhood.

\subsubsection{5. Geometric meaning of \(A_{\mathfrak p}\)}

The local ring
\[
A_{\mathfrak p}
\]
captures what happens \emph{near} the point \(\mathfrak p\).

Globally, the ring \(A\) contains information about the whole space \(\Spec(A)\). After passing to
\[
A_{\mathfrak p},
\]
we invert all elements not in \(\mathfrak p\), which means that anything nonvanishing at \(\mathfrak p\) becomes invertible. Thus only the behavior near \(\mathfrak p\) remains visible.

This is why \(A_{\mathfrak p}\) is the algebraic analogue of looking at germs of functions near a point in analysis or differential geometry.

\subsubsection{6. Geometric meaning of \(\kappa(\mathfrak p)\)}

The residue field
\[
\kappa(\mathfrak p)=A_{\mathfrak p}/\mathfrak p A_{\mathfrak p}
\]
captures the \emph{value} of a function at the point \(\mathfrak p\).

If
\[
\frac{a}{s}\in A_{\mathfrak p},
\]
its image in \(\kappa(\mathfrak p)\) is the ``value'' of that germ at the point \(\mathfrak p\). Since \(s\notin \mathfrak p\), its image is nonzero in \(A/\mathfrak p\), hence invertible in \(\kappa(\mathfrak p)\), so evaluation makes sense.

Thus the map
\[
A_{\mathfrak p}\longrightarrow \kappa(\mathfrak p)
\]
plays the role of evaluating a regular function at the point \(\mathfrak p\).

When \(\kappa(\mathfrak p)=k\), one may think of the point as a \(k\)-rational point. When \(\kappa(\mathfrak p)\) is a larger field extension of \(k\), the point is more general.

\subsubsection{7. Example: \(A=\mathbb Z\)}

Take
\[
A=\mathbb Z.
\]
Its prime ideals are
\[
(0),\ (2),\ (3),\ (5),\dots
\]

\paragraph{The point \(\mathfrak p=(5)\).}

Then
\[
A_{\mathfrak p}=\mathbb Z_{(5)}
=
\left\{
\frac{a}{b}\in \mathbb Q:5\nmid b
\right\}.
\]
Its maximal ideal is
\[
(5)\mathbb Z_{(5)},
\]
and the residue field is
\[
\kappa((5))
=
\mathbb Z_{(5)}/5\mathbb Z_{(5)}
\cong
\mathbb F_5.
\]

Geometrically, the point \((5)\) is a closed point of \(\Spec(\mathbb Z)\), and the residue field tells us that reducing modulo this point gives values in
\[
\mathbb F_5.
\]

\paragraph{The point \(\mathfrak p=(0)\).}

Then
\[
A_{(0)}=\mathbb Q,
\]
because in a domain localization at \((0)\) gives the field of fractions. Also
\[
\kappa((0))=\mathbb Q.
\]

This is the generic point of \(\Spec(\mathbb Z)\). It is not closed, and it sees the whole space in a generic way.

\subsubsection{8. Example: \(A=k[x]\)}

Let
\[
A=k[x].
\]

\paragraph{The point \(\mathfrak p=(x-a)\), where \(a\in k\).}

Then
\[
A_{\mathfrak p}=k[x]_{(x-a)},
\]
the local ring of functions regular near \(x=a\). Its maximal ideal is
\[
(x-a)k[x]_{(x-a)},
\]
and the residue field is
\[
\kappa((x-a))
\cong
k[x]/(x-a)
\cong
k.
\]

So at a \(k\)-rational point \(a\), the value of a regular function lies in \(k\).

\paragraph{The point \(\mathfrak p=(f)\), where \(f\in k[x]\) is irreducible.}

Then
\[
A_{\mathfrak p}=k[x]_{(f)},
\]
and
\[
\kappa((f))
\cong
k[x]/(f).
\]

If \(f\) has degree \(>1\), then \(k[x]/(f)\) is a field extension of \(k\). So the point is not necessarily \(k\)-rational, even though it is still a point of \(\Spec(k[x])\).

\subsubsection{9. Example: \(A=k[x,y]\)}

Let
\[
A=k[x,y].
\]

\paragraph{The point \(\mathfrak p=(x,y)\).}

Then
\[
A_{\mathfrak p}=k[x,y]_{(x,y)},
\]
which is the local ring at the origin. Its maximal ideal is
\[
(x,y)A_{\mathfrak p},
\]
and the residue field is
\[
\kappa((x,y))
\cong
k.
\]

Geometrically, this is the point \((0,0)\) in the affine plane.

\paragraph{The generic point \(\mathfrak p=(0)\).}

Then
\[
A_{(0)}=\operatorname{Frac}(k[x,y])=k(x,y),
\]
the field of rational functions in two variables, and
\[
\kappa((0))=k(x,y).
\]

This is the generic point of the whole affine plane.

\paragraph{A non-closed prime ideal, such as \(\mathfrak p=(y)\).}

Then
\[
A_{\mathfrak p}=k[x,y]_{(y)},
\]
and
\[
\kappa((y))
\cong
\operatorname{Frac}(k[x,y]/(y))
\cong
\operatorname{Frac}(k[x))
\cong
k(x).
\]

Geometrically, \((y)\) corresponds not to a single closed point, but to the generic point of the line
\[
V(y).
\]
Its residue field is the function field of that line.

\subsubsection{10. Closed points and generic points}

In \(\Spec(A)\), not every point is closed.

\paragraph{Closed points.}
A point \(\mathfrak p\) is closed if and only if \(\mathfrak p\) is maximal. In that case,
\[
\kappa(\mathfrak p)=A/\mathfrak p.
\]

\paragraph{Generic points.}
A non-maximal prime ideal gives a non-closed point. Such a point is often a generic point of an irreducible closed subset. Its residue field is then typically a field of rational functions rather than just the base field.

For example, in
\[
\Spec(k[x,y]),
\]
the prime ideal
\[
(y)
\]
is the generic point of the line \(y=0\), and
\[
\kappa((y))\cong k(x).
\]

\subsubsection{11. Functions vanishing at a point}

The maximal ideal
\[
\mathfrak p A_{\mathfrak p}
\]
consists of germs of functions vanishing at the point \(\mathfrak p\). Passing to the quotient
\[
A_{\mathfrak p}/\mathfrak p A_{\mathfrak p}
\]
kills all such functions and keeps only their value at the point.

This is exactly analogous to the situation in differential geometry where germs of smooth functions at a point modulo those vanishing at the point give the ground field.

\subsubsection{12. Summary}

For a point
\[
\mathfrak p\in \Spec(A),
\]
the local ring
\[
A_{\mathfrak p}
\]
is the stalk of the structure sheaf at \(\mathfrak p\), and it records the behavior of regular functions near that point.

Its maximal ideal is
\[
\mathfrak p A_{\mathfrak p},
\]
and the quotient
\[
\kappa(\mathfrak p)=A_{\mathfrak p}/\mathfrak p A_{\mathfrak p}
\]
is the residue field, which records the value of local functions at the point.

One has the exact sequence
\[
0
\longrightarrow
\mathfrak p A_{\mathfrak p}
\longrightarrow
A_{\mathfrak p}
\longrightarrow
\kappa(\mathfrak p)
\longrightarrow
0,
\]
and also the identification
\[
\kappa(\mathfrak p)\cong \operatorname{Frac}(A/\mathfrak p).
\]

Geometrically:

\[
A_{\mathfrak p}
\quad\text{means ``functions near the point''},
\]
\[
\mathfrak p A_{\mathfrak p}
\quad\text{means ``functions vanishing at the point''},
\]
\[
\kappa(\mathfrak p)
\quad\text{means ``the field of values at the point.''}
\]

\subsection{Examples and counterexamples of free, flat, and faithful modules}

In this subsection we collect concrete examples of modules that are free, flat, and faithful, together with important counterexamples. Throughout, \(A\) is a commutative ring with \(1\), and all modules are unital \(A\)-modules.

\subsubsection{1. Definitions}

An \(A\)-module \(M\) is called \emph{free} if there exists a set \(I\) such that
\[
M \cong \bigoplus_{i\in I} A.
\]
Equivalently, \(M\) has a basis over \(A\).

An \(A\)-module \(M\) is called \emph{flat} if for every injective homomorphism of \(A\)-modules
\[
N' \hookrightarrow N,
\]
the induced map
\[
N'\otimes_A M \longrightarrow N\otimes_A M
\]
is injective. Equivalently, tensoring with \(M\) preserves short exact sequences.

An \(A\)-module \(M\) is called \emph{faithful} if
\[
\Ann_A(M)=0.
\]
That is, the only element of \(A\) acting as zero on all of \(M\) is \(0\).

\subsubsection{2. First relations between the notions}

Every free module is flat.

Indeed, if
\[
M\cong \bigoplus_{i\in I} A,
\]
then
\[
N\otimes_A M
\cong
\bigoplus_{i\in I} (N\otimes_A A)
\cong
\bigoplus_{i\in I} N,
\]
so tensoring with \(M\) preserves exactness.

Every nonzero free module is also faithful.

Indeed, if
\[
M\cong \bigoplus_{i\in I} A
\]
with \(I\neq \varnothing\), and if \(a\in A\) kills all of \(M\), then in particular
\[
a\cdot 1=0
\]
in each copy of \(A\), hence \(a=0\).

Thus:
\[
\text{free} \Longrightarrow \text{flat},
\qquad
\text{nonzero free} \Longrightarrow \text{faithful}.
\]

The converses are false in general.

\subsubsection{3. Free modules: basic examples}

\paragraph{Example 1: \(A\) itself.}

The ring \(A\), viewed as an \(A\)-module over itself, is free of rank \(1\):
\[
A \cong A.
\]

\paragraph{Example 2: finite-rank free modules.}

For any integer \(n\geq 1\),
\[
A^n=A\oplus \cdots \oplus A
\]
is free of rank \(n\).

Over \(\mathbb Z\), this is the same as a free abelian group:
\[
\mathbb Z^n.
\]

\paragraph{Example 3: polynomial rings.}

Over
\[
A=k[x],
\]
the module
\[
A^m = k[x]^m
\]
is free of rank \(m\).

\paragraph{Example 4: principal ideals generated by a non-zero-divisor.}

Let
\[
A=k[x].
\]
Then the ideal
\[
(x)\subseteq k[x]
\]
is free of rank \(1\) as an \(A\)-module, because multiplication by \(x\) gives an isomorphism
\[
A \longrightarrow (x),
\qquad
f \longmapsto xf.
\]

More generally, if \(a\in A\) is not a zero-divisor and the map
\[
A \longrightarrow (a),\qquad r\longmapsto ra
\]
is injective and surjective, then \((a)\cong A\) as \(A\)-modules, hence \((a)\) is free of rank \(1\).

\subsubsection{4. Modules that are not free}

\paragraph{Example 5: \(\mathbb Z/n\mathbb Z\) over \(\mathbb Z\).}

The module
\[
\mathbb Z/n\mathbb Z
\]
is not free as a \(\mathbb Z\)-module for \(n>1\), because every free abelian group is torsion-free, whereas
\[
n\cdot \overline{1}=0
\]
in \(\mathbb Z/n\mathbb Z\).

\paragraph{Example 6: \(\mathbb Q\) over \(\mathbb Z\).}

The \(\mathbb Z\)-module
\[
\mathbb Q
\]
is not free. Indeed, a free abelian group has a basis and therefore is a direct sum of copies of \(\mathbb Z\), but \(\mathbb Q\) is divisible:
\[
\forall q\in \mathbb Q,\ \forall n\ge 1,\ \exists r\in \mathbb Q \text{ with } nr=q.
\]
A nonzero free abelian group is never divisible.

\paragraph{Example 7: localization \(\mathbb Z_{(p)}\) over \(\mathbb Z\).}

The module
\[
\mathbb Z_{(p)}
=
\left\{
\frac{a}{b}\in \mathbb Q : p\nmid b
\right\}
\]
is not free as a \(\mathbb Z\)-module. It is torsion-free, but not free.

A quick reason is that it is not cyclic:
\[
1 \text{ does not generate } \frac{1}{q}
\quad\text{for } q\neq p.
\]
But it also cannot be a free abelian group of rank \(>1\), since it embeds into \(\mathbb Q\), and subgroups of \(\mathbb Q\) have very special structure. In particular, \(\mathbb Z_{(p)}\) is a standard example of a flat module that is not free.

\paragraph{Example 8: residue fields are usually not free.}

Let
\[
A=k[x],
\qquad
M=A/(x)\cong k.
\]
As an \(A\)-module, \(k\) is not free, because it has torsion with respect to \(x\):
\[
x\cdot \overline{1}=0.
\]
So quotient modules like \(A/I\) are usually not free unless \(I=0\).

\subsubsection{5. Flat modules: main examples}

\paragraph{Example 9: all free modules are flat.}

As noted above, every module of the form
\[
\bigoplus_{i\in I} A
\]
is flat.

So in particular:
\[
A,\quad A^n,\quad A^{(I)}
\]
are flat.

\paragraph{Example 10: localizations are flat.}

For every multiplicative set \(S\subseteq A\), the localization
\[
S^{-1}A
\]
is flat over \(A\).

In particular, for every prime ideal \(\mathfrak p\),
\[
A_{\mathfrak p}
\]
is flat over \(A\).

Examples:
\[
\mathbb Z_{(p)} \text{ is flat over } \mathbb Z,
\]
\[
k[x]_{(x)} \text{ is flat over } k[x],
\]
\[
k[x,y]_{(x,y)} \text{ is flat over } k[x,y].
\]

\paragraph{Example 11: \(\mathbb Q\) is flat over \(\mathbb Z\).}

Since
\[
\mathbb Q = \mathbb Z_{(0)}
\]
is the localization of \(\mathbb Z\) at all nonzero integers, it is flat over \(\mathbb Z\).

\paragraph{Example 12: vector spaces over a field.}

If \(A=k\) is a field, then every \(k\)-module is a vector space, hence free, hence flat.

So over a field:
\[
\text{all modules are free and flat.}
\]

\paragraph{Example 13: over a PID, torsion-free implies flat.}

Over a principal ideal domain such as
\[
\mathbb Z
\quad\text{or}\quad
k[x],
\]
every torsion-free module is flat.

Thus:
\[
\mathbb Q,\quad \mathbb Z_{(p)},\quad (x)\subseteq k[x]
\]
are flat.

This is an important source of examples of flat modules that are not free.

\subsubsection{6. Modules that are not flat}

\paragraph{Example 14: \(\mathbb Z/n\mathbb Z\) over \(\mathbb Z\).}

The module
\[
\mathbb Z/n\mathbb Z
\]
is not flat over \(\mathbb Z\) for \(n>1\).

Indeed, consider the exact sequence
\[
0
\longrightarrow
\mathbb Z
\overset{\cdot n}{\longrightarrow}
\mathbb Z
\longrightarrow
\mathbb Z/n\mathbb Z
\longrightarrow
0.
\]
Tensoring with \(\mathbb Z/n\mathbb Z\) gives
\[
0
\longrightarrow
\mathbb Z/n\mathbb Z
\overset{\cdot n}{\longrightarrow}
\mathbb Z/n\mathbb Z
\longrightarrow
(\mathbb Z/n\mathbb Z)\otimes_{\mathbb Z}(\mathbb Z/n\mathbb Z)
\longrightarrow
0.
\]
But multiplication by \(n\) on \(\mathbb Z/n\mathbb Z\) is the zero map, so the first map is not injective. Hence \(\mathbb Z/n\mathbb Z\) is not flat.

\paragraph{Example 15: residue field \(k\) over \(k[x]\).}

Let
\[
A=k[x],
\qquad
M=A/(x)\cong k.
\]
Consider the exact sequence
\[
0
\longrightarrow
A
\overset{\cdot x}{\longrightarrow}
A
\longrightarrow
A/(x)
\longrightarrow
0.
\]
Tensoring with \(M\) gives
\[
0
\longrightarrow
M
\overset{\cdot x}{\longrightarrow}
M
\longrightarrow
M\otimes_A A/(x)
\longrightarrow
0.
\]
But multiplication by \(x\) on \(M=A/(x)\) is zero, so again the first map is not injective. Therefore \(k\) is not flat over \(k[x]\).

\paragraph{Example 16: quotient modules are usually not flat.}

In general,
\[
A/I
\]
is usually not flat over \(A\) unless \(I\) is very special. The previous examples
\[
\mathbb Z/n\mathbb Z,\qquad k[x]/(x)
\]
show this clearly.

\subsubsection{7. Faithful modules: main examples}

\paragraph{Example 17: \(A\) over itself.}

The module \(A\) is faithful over \(A\), because
\[
\Ann_A(A)=0.
\]

\paragraph{Example 18: every nonzero free module is faithful.}

If
\[
M\cong A^n
\qquad (n\ge 1),
\]
then
\[
\Ann_A(M)=0.
\]
So
\[
A^n
\]
is faithful.

\paragraph{Example 19: \(\mathbb Q\) over \(\mathbb Z\).}

The \(\mathbb Z\)-module \(\mathbb Q\) is faithful, because if an integer \(m\) kills all rational numbers, then in particular
\[
m\cdot 1=0
\]
in \(\mathbb Q\), hence \(m=0\).

Thus
\[
\Ann_{\mathbb Z}(\mathbb Q)=0.
\]

\paragraph{Example 20: localizations over domains.}

If \(A\) is a domain, then every localization \(S^{-1}A\) is faithful as an \(A\)-module.

Indeed, if
\[
a\in A
\]
kills all of \(S^{-1}A\), then in particular
\[
a\cdot \frac{1}{1}=0,
\]
so
\[
\frac{a}{1}=0
\]
in \(S^{-1}A\). Since \(A\) is a domain, this implies \(a=0\).

Examples:
\[
\mathbb Z_{(p)} \text{ is faithful over } \mathbb Z,
\]
\[
k[x]_{(x)} \text{ is faithful over } k[x].
\]

\paragraph{Example 21: the field of fractions of a domain.}

If \(A\) is a domain with field of fractions
\[
K=\Frac(A),
\]
then \(K\) is a faithful \(A\)-module.

For example:
\[
\mathbb Q \text{ is faithful over } \mathbb Z,
\]
\[
k(x) \text{ is faithful over } k[x].
\]

\subsubsection{8. Modules that are not faithful}

\paragraph{Example 22: quotient modules \(A/I\).}

For any ideal \(I\subseteq A\),
\[
\Ann_A(A/I)=I.
\]
Indeed, \(a\in A\) kills \(A/I\) if and only if
\[
a\cdot (1+I)=0
\]
in \(A/I\), which means
\[
a\in I.
\]
Therefore \(A/I\) is faithful if and only if
\[
I=0.
\]

So:
\[
\mathbb Z/n\mathbb Z
\]
is not faithful as a \(\mathbb Z\)-module for \(n>1\), and
\[
k[x]/(x)\cong k
\]
is not faithful as a \(k[x]\)-module.

\paragraph{Example 23: residue fields are not faithful in general.}

Let
\[
A=k[x],
\qquad
M=A/(x)\cong k.
\]
Then
\[
x\cdot M=0,
\]
so
\[
\Ann_A(M)\supseteq (x)\neq 0.
\]
Hence \(k\) is not faithful as an \(A\)-module.

\paragraph{Example 24: a module over a product ring.}

Let
\[
A=k\times k,
\]
and consider the \(A\)-module
\[
M=k\times 0.
\]
Then
\[
(0,1)\cdot (a,0)=(0,0)
\]
for all \((a,0)\in M\), so
\[
(0,1)\in \Ann_A(M).
\]
Hence \(M\) is not faithful.

This is a useful example showing that even very natural modules over a ring need not be faithful.

\subsubsection{9. A module can be flat and faithful but not free}

This is a very important phenomenon.

\paragraph{Example 25: \(\mathbb Q\) over \(\mathbb Z\).}

The module
\[
\mathbb Q
\]
is flat over \(\mathbb Z\), because it is a localization. It is faithful over \(\mathbb Z\), because no nonzero integer kills all of \(\mathbb Q\). But it is not free as a \(\mathbb Z\)-module.

So:
\[
\mathbb Q \text{ is flat and faithful, but not free.}
\]

\paragraph{Example 26: \(\mathbb Z_{(p)}\) over \(\mathbb Z\).}

Likewise,
\[
\mathbb Z_{(p)}
\]
is flat and faithful over \(\mathbb Z\), but not free.

\paragraph{Example 27: \(k[x]_{(x)}\) over \(k[x]\).}

The module
\[
k[x]_{(x)}
\]
is flat and faithful over \(k[x]\), but not free.

These examples are fundamental in commutative algebra.

\subsubsection{10. A module can be faithful but not flat}

\paragraph{Example 28: ideals in rings with zero divisors.}

Let
\[
A=k[\varepsilon]/(\varepsilon^2),
\qquad
M=(\varepsilon).
\]
Then \(M\cong k\) as an abelian group, and
\[
\varepsilon\cdot M=0.
\]
So \(M\) is not faithful.

Thus this example does not yet separate faithfulness from flatness.

A better source of examples comes from more complicated rings, but in ordinary introductory examples the main phenomenon is usually:
\[
\text{faithful does not imply flat,}
\]
although explicit elementary examples are less canonical than the standard localization examples above.

\subsubsection{11. A module can be flat but not faithful}

\paragraph{Example 29: localizations over rings with zero divisors.}

Let
\[
A=\mathbb Z/6\mathbb Z,
\qquad
\mathfrak p=(2).
\]
Then
\[
A_{\mathfrak p}
\]
is flat over \(A\), because every localization is flat. But it is not faithful, because some nonzero elements of \(A\) become zero in the localization.

Thus:
\[
\text{flat does not imply faithful.}
\]

\subsubsection{12. A module can be faithful but not free}

\paragraph{Example 30: \(\mathbb Q\) over \(\mathbb Z\).}

As above,
\[
\mathbb Q
\]
is faithful but not free.

\paragraph{Example 31: \(A_{\mathfrak p}\) over a domain \(A\).}

If \(A\) is a domain, then \(A_{\mathfrak p}\) is faithful, and usually not free.

So:
\[
\text{faithful does not imply free.}
\]

\subsubsection{13. Summary table of examples}

\[
\begin{array}{|c|c|c|c|}
	\hline
	\text{Module} & \text{Free} & \text{Flat} & \text{Faithful} \\
	\hline
	A & \text{yes} & \text{yes} & \text{yes} \\
	\hline
	A^n & \text{yes} & \text{yes} & \text{yes if } n\ge 1 \\
	\hline
	\mathbb Z/n\mathbb Z \text{ over } \mathbb Z & \text{no} & \text{no} & \text{no} \\
	\hline
	\mathbb Q \text{ over } \mathbb Z & \text{no} & \text{yes} & \text{yes} \\
	\hline
	\mathbb Z_{(p)} \text{ over } \mathbb Z & \text{no} & \text{yes} & \text{yes} \\
	\hline
	k[x]/(x) \text{ over } k[x] & \text{no} & \text{no} & \text{no} \\
	\hline
	k[x]_{(x)} \text{ over } k[x] & \text{no in general} & \text{yes} & \text{yes} \\
	\hline
	A/I \text{ over } A,\ I\neq 0 & \text{usually no} & \text{usually no} & \text{no} \\
	\hline
\end{array}
\]

\subsubsection{14. Final guiding picture}

One should remember the following standard examples:

\[
A,\ A^n
\quad\text{are free, flat, and faithful;}
\]

\[
\mathbb Z/n\mathbb Z,\ A/I
\quad\text{are usually neither free nor flat nor faithful;}
\]

\[
\mathbb Q,\ \mathbb Z_{(p)},\ A_{\mathfrak p}
\quad\text{are standard examples of modules that are flat, often faithful, but not free.}
\]

In particular, localization provides one of the most important classes of flat modules in commutative algebra, and quotient modules provide one of the most important classes of non-flat and non-faithful modules.

\subsection{Concrete examples of associated graded rings \(\operatorname{gr}_I(A)\)}

Let \(A\) be a ring and \(I\subseteq A\) an ideal. Recall that
\[
\operatorname{gr}_I(A)
=
\bigoplus_{n\ge 0} I^n/I^{n+1}.
\]
Multiplication is induced from multiplication in \(A\):
\[
(a+I^{r+1})(b+I^{s+1})=ab+I^{r+s+1}
\]
for
\[
a\in I^r,\qquad b\in I^s.
\]

If
\[
a\in I^n\setminus I^{n+1},
\]
its image in
\[
I^n/I^{n+1}
\]
is often denoted by
\[
a^*.
\]
This is called the \emph{initial form} of \(a\).

The following examples make the construction concrete.

\subsubsection{1. The case \(A=k[x]\), \(I=(x)\)}

Let
\[
A=k[x],
\qquad
I=(x).
\]
Then
\[
I^n=(x^n)
\]
for every \(n\ge 0\). Hence
\[
I^n/I^{n+1}=(x^n)/(x^{n+1}).
\]
Each quotient
\[
(x^n)/(x^{n+1})
\]
is a one-dimensional \(k\)-vector space generated by the class of \(x^n\).

Let
\[
t:=x+(x^2)\in I/I^2.
\]
Then
\[
t^n=x^n+(x^{n+1})\in I^n/I^{n+1}.
\]
So every graded piece is generated by \(t^n\), and the multiplication behaves exactly as in a polynomial ring. Therefore
\[
\operatorname{gr}_{(x)}(k[x])\cong k[t].
\]

Thus in this example:

\[
\operatorname{gr}_{(x)}(k[x]) \text{ is Noetherian,}
\]
\[
\operatorname{gr}_{(x)}(k[x]) \text{ is an integral domain.}
\]

\subsubsection{2. The local case \(A=k[x]_{(x)}\), \(\mathfrak m=(x)\)}

Now let
\[
A=k[x]_{(x)},
\qquad
\mathfrak m=(x).
\]
Then
\[
\mathfrak m^n=(x^n),
\]
and the same computation as above gives
\[
\mathfrak m^n/\mathfrak m^{n+1}\cong k\cdot x^n.
\]
Hence
\[
\operatorname{gr}_{\mathfrak m}(A)\cong k[t].
\]

This is a concrete illustration of part (b) in the problem: since
\[
\operatorname{gr}_{\mathfrak m}(A)
\]
is a domain, the local ring
\[
A=k[x]_{(x)}
\]
is a domain.

\subsubsection{3. The case \(A=\mathbb Z\), \(I=(p)\)}

Let
\[
A=\mathbb Z,
\qquad
I=(p),
\]
where \(p\) is a prime number. Then
\[
I^n=(p^n).
\]
So
\[
I^n/I^{n+1}=p^n\mathbb Z/p^{n+1}\mathbb Z.
\]
Each graded piece is a one-dimensional vector space over
\[
\mathbb Z/p\mathbb Z=\mathbb F_p,
\]
generated by the class of \(p^n\).

Let
\[
t:=p+(p^2)\in (p)/(p^2).
\]
Then
\[
t^n=p^n+(p^{n+1}).
\]
Therefore
\[
\operatorname{gr}_{(p)}(\mathbb Z)\cong \mathbb F_p[t].
\]

So again:

\[
\operatorname{gr}_{(p)}(\mathbb Z) \text{ is Noetherian,}
\]
\[
\operatorname{gr}_{(p)}(\mathbb Z) \text{ is a domain.}
\]

If instead
\[
A=\mathbb Z_{(p)},
\qquad
\mathfrak m=(p),
\]
then similarly
\[
\operatorname{gr}_{\mathfrak m}(A)\cong \mathbb F_p[t].
\]

\subsubsection{4. The case \(A=k[x,y]\), \(I=(x,y)\)}

Let
\[
A=k[x,y],
\qquad
I=(x,y).
\]
Then
\[
I^n
\]
consists of all polynomials with no terms of total degree \(<n\). Therefore
\[
I^n/I^{n+1}
\]
is precisely the space of homogeneous polynomials of total degree \(n\).

For example,
\[
I/I^2
=
\langle x^*,y^*\rangle_k,
\]
\[
I^2/I^3
=
\langle (x^*)^2,x^*y^*,(y^*)^2\rangle_k,
\]
\[
I^3/I^4
=
\langle (x^*)^3,(x^*)^2y^*,x^*(y^*)^2,(y^*)^3\rangle_k,
\]
and so on.

If we introduce formal variables
\[
X,\ Y
\]
and send
\[
X\longmapsto x^*=x+I^2,
\qquad
Y\longmapsto y^*=y+I^2,
\]
then every homogeneous polynomial in \(X,Y\) of degree \(n\) corresponds to an element of
\[
I^n/I^{n+1}.
\]
This gives an isomorphism
\[
\operatorname{gr}_{(x,y)}(k[x,y])\cong k[X,Y].
\]

Thus
\[
\operatorname{gr}_{(x,y)}(k[x,y])
\]
is a polynomial ring in two variables over \(k\), hence is Noetherian and a domain.

\subsubsection{5. The local two-variable case \(A=k[x,y]_{(x,y)}\)}

Let
\[
A=k[x,y]_{(x,y)},
\qquad
\mathfrak m=(x,y).
\]
The same reasoning as above shows that
\[
\mathfrak m^n/\mathfrak m^{n+1}
\]
is the space of homogeneous forms of total degree \(n\), so
\[
\operatorname{gr}_{\mathfrak m}(A)\cong k[X,Y].
\]

Thus
\[
\operatorname{gr}_{\mathfrak m}(A)
\]
is a domain, and therefore
\[
A=k[x,y]_{(x,y)}
\]
is a domain, as expected.

\subsubsection{6. The power series case \(A=k[[x,y]]\), \(\mathfrak m=(x,y)\)}

Let
\[
A=k[[x,y]],
\qquad
\mathfrak m=(x,y).
\]
Then
\[
\mathfrak m^n
\]
consists of power series all of whose monomials have total degree at least \(n\). Hence
\[
\mathfrak m^n/\mathfrak m^{n+1}
\]
is again the space of homogeneous polynomials of total degree \(n\).

Therefore
\[
\operatorname{gr}_{\mathfrak m}(k[[x,y]])\cong k[X,Y].
\]

So the associated graded ring is Noetherian and a domain.

More generally, for
\[
A=k[[x_1,\dots,x_r]],
\qquad
\mathfrak m=(x_1,\dots,x_r),
\]
one has
\[
\operatorname{gr}_{\mathfrak m}(A)\cong k[X_1,\dots,X_r].
\]

\subsubsection{7. A non-domain example: dual numbers}

Let
\[
A=k[\varepsilon]/(\varepsilon^2),
\qquad
\mathfrak m=(\varepsilon).
\]
Then
\[
\mathfrak m^2=0.
\]
So
\[
\operatorname{gr}_{\mathfrak m}(A)
=
A/\mathfrak m \oplus \mathfrak m/\mathfrak m^2.
\]
Now
\[
A/\mathfrak m \cong k,
\qquad
\mathfrak m/\mathfrak m^2 \cong k\varepsilon^*.
\]
Thus
\[
\operatorname{gr}_{\mathfrak m}(A)\cong k[T]/(T^2),
\]
where
\[
T=\varepsilon^*.
\]

This is \emph{not} a domain, because
\[
T^2=0.
\]
And indeed the original ring
\[
A=k[\varepsilon]/(\varepsilon^2)
\]
is not a domain, since
\[
\varepsilon^2=0.
\]

This is a very concrete example showing how zero divisors and nilpotents in \(A\) can appear in
\[
\operatorname{gr}_{\mathfrak m}(A).
\]

\subsubsection{8. A reducible singular example: \(A=k[[x,y]]/(xy)\)}

Let
\[
A=k[[x,y]]/(xy),
\qquad
\mathfrak m=(x,y)/(xy).
\]
Then
\[
x^*=x+\mathfrak m^2,
\qquad
y^*=y+\mathfrak m^2
\]
are nonzero elements of
\[
\mathfrak m/\mathfrak m^2.
\]
But since
\[
xy=0
\]
already in \(A\), we get
\[
x^*y^*=0
\]
in
\[
\operatorname{gr}_{\mathfrak m}(A).
\]
In fact,
\[
\operatorname{gr}_{\mathfrak m}(A)\cong k[X,Y]/(XY).
\]

So
\[
\operatorname{gr}_{\mathfrak m}(A)
\]
is not a domain, and indeed
\[
A
\]
itself is not a domain.

\subsubsection{9. A warning example: \(A\) may be a domain while \(\operatorname{gr}_{\mathfrak m}(A)\) is not}

The implication in part (b) of the problem is
\[
\operatorname{gr}_{\mathfrak m}(A)\text{ is a domain } \Longrightarrow A\text{ is a domain.}
\]
The converse is false in general.

Consider
\[
A=k[[t^2,t^3]]\subseteq k[[t]].
\]
This is a domain because it is a subring of the domain
\[
k[[t]].
\]
Let
\[
\mathfrak m=(t^2,t^3).
\]
Then
\[
t^3\notin \mathfrak m^2,
\]
so
\[
y^*:=t^3+\mathfrak m^2
\]
is a nonzero element of
\[
\mathfrak m/\mathfrak m^2.
\]
But
\[
(y^*)^2=t^6+\mathfrak m^3=0
\]
in
\[
\mathfrak m^2/\mathfrak m^3,
\]
because
\[
t^6\in \mathfrak m^3.
\]
Hence
\[
\operatorname{gr}_{\mathfrak m}(A)
\]
has a nonzero nilpotent element, so it is not a domain.

Thus \(A\) can be a domain while
\[
\operatorname{gr}_{\mathfrak m}(A)
\]
fails to be one.

\subsubsection{10. General pattern in regular local rings}

These examples suggest an important general pattern. If
\[
(A,\mathfrak m)
\]
is a regular local ring and
\[
\mathfrak m=(x_1,\dots,x_r),
\]
then one often has
\[
\operatorname{gr}_{\mathfrak m}(A)\cong (A/\mathfrak m)[X_1,\dots,X_r].
\]

This explains why in the standard regular examples one gets polynomial rings:

\[
\operatorname{gr}_{(x)}(k[x])\cong k[t],
\]
\[
\operatorname{gr}_{(p)}(\mathbb Z)\cong \mathbb F_p[t],
\]
\[
\operatorname{gr}_{(x,y)}(k[x,y])\cong k[X,Y],
\]
\[
\operatorname{gr}_{(x,y)}(k[[x,y]])\cong k[X,Y].
\]

So in many familiar regular situations, the associated graded ring is especially simple: it is just a polynomial ring over the residue field.

\subsubsection{11. Summary of the main examples}

\paragraph{Domain examples.}
\[
\operatorname{gr}_{(x)}(k[x])\cong k[t],
\]
\[
\operatorname{gr}_{(p)}(\mathbb Z)\cong \mathbb F_p[t],
\]
\[
\operatorname{gr}_{(x,y)}(k[x,y])\cong k[X,Y],
\]
\[
\operatorname{gr}_{(x,y)}(k[[x,y]])\cong k[X,Y].
\]

All of these are Noetherian domains.

\paragraph{Non-domain examples.}
\[
\operatorname{gr}_{(\varepsilon)}\bigl(k[\varepsilon]/(\varepsilon^2)\bigr)\cong k[T]/(T^2),
\]
\[
\operatorname{gr}_{\mathfrak m}\bigl(k[[x,y]]/(xy)\bigr)\cong k[X,Y]/(XY).
\]

These are not domains.

\paragraph{Warning example.}
\[
A=k[[t^2,t^3]]
\]
is a domain, but
\[
\operatorname{gr}_{\mathfrak m}(A)
\]
is not a domain.

So the condition that
\[
\operatorname{gr}_{\mathfrak m}(A)
\]
be a domain is strong enough to force \(A\) to be a domain, but it is not necessary.

\subsubsection{Comparison of \(k[x]_{(x)}\) and \(k[x]/(x)\)}

The rings
\[
k[x]_{(x)}
\qquad\text{and}\qquad
k[x]/(x)
\]
are \emph{not} the same.

They arise from two different constructions:

\paragraph{1. The quotient ring \(k[x]/(x)\).}

This is obtained by factoring out the ideal
\[
(x).
\]
In other words, one forces
\[
x=0.
\]
Every polynomial is identified with its constant term, so there is a natural isomorphism
\[
k[x]/(x)\cong k.
\]

Explicitly, the quotient map is
\[
k[x]\longrightarrow k[x]/(x),
\qquad
f(x)\longmapsto f(x)+(x),
\]
and under the identification with \(k\), this is just evaluation at \(x=0\):
\[
f(x)\longmapsto f(0).
\]

\paragraph{2. The localization \(k[x]_{(x)}\).}

This is obtained by inverting all elements not belonging to the prime ideal
\[
(x).
\]
So
\[
k[x]_{(x)}=S^{-1}k[x],
\qquad
S=k[x]\setminus (x).
\]
Since a polynomial lies outside \((x)\) exactly when its constant term is nonzero, the elements of
\[
k[x]_{(x)}
\]
are fractions
\[
\frac{f(x)}{g(x)}
\qquad\text{with}\qquad g(0)\neq 0.
\]

Thus in \(k[x]_{(x)}\), the element \(x\) does \emph{not} become zero. Instead, \(x\) remains a nonunit and belongs to the maximal ideal
\[
(x)k[x]_{(x)}.
\]

\paragraph{3. The essential difference.}

In the quotient ring
\[
k[x]/(x),
\]
one has
\[
x=0.
\]

In the localization
\[
k[x]_{(x)},
\]
one does \emph{not} kill \(x\). One only inverts those polynomials that do not vanish at \(x=0\).

So:
\[
k[x]/(x)\cong k,
\]
whereas
\[
k[x]_{(x)}
\]
is a larger local ring consisting of rational expressions regular at \(x=0\).

Therefore
\[
k[x]_{(x)}\neq k[x]/(x).
\]

\paragraph{4. The relation between them.}

Although they are not equal, they are closely related. The ring
\[
k[x]_{(x)}
\]
has maximal ideal
\[
(x)k[x]_{(x)},
\]
and its residue field is
\[
k[x]_{(x)}/(x)k[x]_{(x)}\cong k.
\]

Hence one has
\[
k[x]_{(x)}/(x)k[x]_{(x)}
\cong
k[x]/(x)
\cong
k.
\]

So the correct relationship is:
\[
k[x]/(x)
\cong
k[x]_{(x)}/(x)k[x]_{(x)},
\]
not
\[
k[x]_{(x)}=k[x]/(x).
\]

\paragraph{5. Geometric interpretation.}

The ring
\[
k[x]/(x)\cong k
\]
describes the point
\[
x=0
\]
itself.

The ring
\[
k[x]_{(x)}
\]
describes functions defined \emph{near} the point \(x=0\).

Thus:
\[
k[x]/(x)
\quad\text{corresponds to the value at the point,}
\]
while
\[
k[x]_{(x)}
\quad\text{corresponds to germs of functions near the point.}
\]

\subsection{Localization in polynomial rings: concrete descriptions and examples}

Let \(k\) be a field. Polynomial rings are among the best places to understand localization concretely, because one can see explicitly which denominators are allowed and what kind of geometric information the localization captures.

\subsubsection{1. General definition}

Let \(A\) be a commutative ring and let \(S\subseteq A\) be a multiplicative set, that is,
\[
1\in S,
\qquad
s,t\in S \Longrightarrow st\in S.
\]
The localization of \(A\) at \(S\) is the ring
\[
S^{-1}A.
\]
Its elements are fractions
\[
\frac{a}{s},
\qquad
a\in A,\ s\in S,
\]
with the usual equivalence relation:
\[
\frac{a}{s}=\frac{b}{t}
\quad\Longleftrightarrow\quad
\exists u\in S \text{ such that } u(at-bs)=0.
\]

The intuitive meaning is simple:

\[
S^{-1}A
\quad\text{is obtained from }A\text{ by forcing every }s\in S\text{ to become invertible.}
\]

In polynomial rings, this means that one starts with ordinary polynomials and then allows division by selected polynomials.

\subsubsection{2. Localization at one polynomial}

Let
\[
A=k[x].
\]
If \(f\in k[x]\), define the multiplicative set
\[
S=\{1,f,f^2,f^3,\dots\}.
\]
Then the localization is written
\[
A_f=k[x]_f.
\]
Its elements are fractions of the form
\[
\frac{g(x)}{f(x)^n},
\qquad
g(x)\in k[x],\ n\ge 0.
\]

\paragraph{Example.}
If
\[
f=x,
\]
then
\[
k[x]_x
\]
consists of fractions
\[
\frac{g(x)}{x^n}.
\]
So
\[
\frac{1}{x}\in k[x]_x,
\qquad
\frac{x^2+1}{x^3}\in k[x]_x,
\qquad
\frac{x+5}{x^2}\in k[x]_x.
\]
But
\[
\frac{1}{x+1}\notin k[x]_x,
\]
because the denominator \(x+1\) is not a power of \(x\).

Thus
\[
k[x]_x
\]
is not the full field of fractions \(k(x)\). It only allows denominators that are powers of \(x\).

\subsubsection{3. Localization at a prime ideal}

Let \(\mathfrak p\subseteq A\) be a prime ideal. Then
\[
A_{\mathfrak p}
\]
means localization at the multiplicative set
\[
S=A\setminus \mathfrak p.
\]
So we invert every element of \(A\) that does not lie in \(\mathfrak p\).

This construction is especially important in algebraic geometry and commutative algebra.

\subsubsection{4. The ring \(k[x]_{(x)}\)}

Take
\[
A=k[x],
\qquad
\mathfrak p=(x).
\]
Then
\[
A_{\mathfrak p}=k[x]_{(x)}.
\]
Since a polynomial lies outside \((x)\) exactly when its constant term is nonzero, the elements of
\[
k[x]_{(x)}
\]
are fractions
\[
\frac{f(x)}{g(x)}
\qquad\text{with}\qquad g(0)\neq 0.
\]

\paragraph{Examples.}
\[
\frac{x^2+1}{x+1}\in k[x]_{(x)},
\]
because
\[
(x+1)(0)=1\neq 0.
\]

\[
\frac{1}{1+x}\in k[x]_{(x)},
\]
because the denominator has nonzero constant term.

But
\[
\frac{1}{x}\notin k[x]_{(x)},
\]
because
\[
x\in (x),
\]
so \(x\) is not inverted in this localization.

Thus \(k[x]_{(x)}\) consists of rational expressions that are regular at the point \(x=0\).

\subsubsection{5. The ring \(k[x]_{(x-a)}\)}

More generally, let
\[
\mathfrak p=(x-a),
\qquad a\in k.
\]
Then
\[
k[x]_{(x-a)}
\]
consists of fractions
\[
\frac{f(x)}{g(x)}
\qquad\text{with}\qquad g(a)\neq 0.
\]

So the denominator must not vanish at \(x=a\).

\paragraph{Examples.}
\[
\frac{x^2+1}{x+1}\in k[x]_{(x-2)},
\]
because
\[
(x+1)(2)=3\neq 0.
\]

But
\[
\frac{1}{x-2}\notin k[x]_{(x-2)},
\]
because the denominator lies in the ideal \((x-2)\).

Geometrically, this ring is the local ring at the point \(x=a\).

\subsubsection{6. The ring \(k[x]_{(0)}\)}

Now let
\[
\mathfrak p=(0).
\]
Since \(k[x]\) is a domain, \((0)\) is a prime ideal. Then
\[
k[x]_{(0)}
\]
means that we invert every nonzero polynomial, because
\[
k[x]\setminus (0)=\{f(x)\in k[x]:f(x)\neq 0\}.
\]

Therefore
\[
k[x]_{(0)}
\]
consists of all fractions
\[
\frac{f(x)}{g(x)}
\qquad\text{with}\qquad g(x)\neq 0.
\]
But this is exactly the field of fractions of \(k[x]\), usually denoted
\[
k(x).
\]
Hence
\[
k[x]_{(0)}=k(x).
\]

So:

\[
k[x]
\subseteq
k[x]_{(x)}
\subseteq
k[x]_{(0)}=k(x).
\]

This chain is extremely useful to keep in mind.

\subsubsection{7. Localization in two variables}

Now let
\[
A=k[x,y].
\]

\paragraph{Localization at \(x\).}
The ring
\[
k[x,y]_x
\]
consists of fractions
\[
\frac{f(x,y)}{x^n}.
\]
Geometrically, this corresponds to the open set where
\[
x\neq 0.
\]

\paragraph{Localization at \((x,y)\).}
The ring
\[
k[x,y]_{(x,y)}
\]
consists of fractions
\[
\frac{f(x,y)}{g(x,y)}
\qquad\text{with}\qquad g(0,0)\neq 0.
\]
Thus one may divide by any polynomial that does not vanish at the origin.

For example,
\[
\frac{x^2+y}{1+x+y}\in k[x,y]_{(x,y)},
\]
because
\[
(1+x+y)(0,0)=1.
\]

But
\[
\frac{1}{x}\notin k[x,y]_{(x,y)},
\qquad
\frac{1}{y}\notin k[x,y]_{(x,y)},
\qquad
\frac{1}{x+y}\notin k[x,y]_{(x,y)},
\]
since all these denominators vanish at \((0,0)\).

Thus
\[
k[x,y]_{(x,y)}
\]
is the local ring of the affine plane at the origin.

\paragraph{Localization at \((y)\).}
The ring
\[
k[x,y]_{(y)}
\]
consists of fractions whose denominator is not divisible by \(y\). For example,
\[
\frac{1}{x+1}\in k[x,y]_{(y)},
\qquad
\frac{1}{x}\in k[x,y]_{(y)},
\]
but
\[
\frac{1}{y}\notin k[x,y]_{(y)}.
\]

Geometrically, the prime ideal \((y)\) corresponds not to a closed point, but to the generic point of the line
\[
y=0.
\]

\subsubsection{8. Comparison with quotients}

It is very important not to confuse localization and quotient.

For example,
\[
k[x]/(x)\cong k,
\]
because in the quotient one forces
\[
x=0.
\]

But
\[
k[x]_{(x)}
\]
does not kill \(x\). Instead, it inverts all polynomials not divisible by \(x\). So
\[
k[x]_{(x)}\neq k[x]/(x).
\]

The correct relationship is
\[
k[x]_{(x)}/(x)k[x]_{(x)}\cong k[x]/(x)\cong k.
\]
Thus the quotient
\[
k[x]_{(x)}/(x)k[x]_{(x)}
\]
is the residue field of the local ring \(k[x]_{(x)}\).

\subsubsection{9. Geometric meaning}

Localization at one polynomial \(f\) corresponds to restricting to the open set where \(f\neq 0\).

Localization at a maximal ideal, such as
\[
(x-a)\subseteq k[x]
\quad\text{or}\quad
(x-a,y-b)\subseteq k[x,y],
\]
corresponds to looking at functions defined near a concrete point.

Localization at a non-maximal prime, such as
\[
(0)\subseteq k[x]
\quad\text{or}\quad
(y)\subseteq k[x,y],
\]
corresponds to looking at functions near a generic point.

So localization can be interpreted as a way to focus algebraically on a chosen part of the space.

\subsubsection{10. Summary}

The most important concrete localizations in polynomial rings are:

\[
k[x]_x
\quad\text{allows denominators }x^n,
\]

\[
k[x]_{(x-a)}
\quad\text{allows denominators not vanishing at }x=a,
\]

\[
k[x]_{(0)}=k(x)
\quad\text{allows all nonzero polynomial denominators,}
\]

\[
k[x,y]_{(x,y)}
\quad\text{allows denominators not vanishing at the origin.}
\]

Thus localization in polynomial rings is the algebraic mechanism for passing from global polynomials to local or partially defined rational expressions.

\section{Worked exercises on membership in localizations}

In the following exercises, the goal is to decide whether a given fraction belongs to a given localization. In each case, the criterion is simply: the denominator must lie in the chosen multiplicative set.

\subsection{Exercises in one variable}

\subsubsection{Exercise 1}

Determine whether
\[
\frac{1}{x}
\]
belongs to
\[
k[x]_x
\qquad\text{and}\qquad
k[x]_{(x)}.
\]

\textbf{Solution.}

In
\[
k[x]_x,
\]
the allowed denominators are powers of \(x\). Since
\[
x=x^1,
\]
we have
\[
\frac{1}{x}\in k[x]_x.
\]

In
\[
k[x]_{(x)},
\]
the allowed denominators are polynomials not in \((x)\), that is, polynomials with nonzero constant term. But
\[
x\in (x),
\]
so
\[
\frac{1}{x}\notin k[x]_{(x)}.
\]

\subsubsection{Exercise 2}

Determine whether
\[
\frac{1}{x+1}
\]
belongs to
\[
k[x]_x
\qquad\text{and}\qquad
k[x]_{(x)}.
\]

\textbf{Solution.}

In
\[
k[x]_x,
\]
the denominator must be a power of \(x\). Since \(x+1\) is not a power of \(x\),
\[
\frac{1}{x+1}\notin k[x]_x.
\]

In
\[
k[x]_{(x)},
\]
the denominator is allowed if it does not lie in \((x)\). Since
\[
(x+1)(0)=1\neq 0,
\]
we have
\[
x+1\notin (x),
\]
hence
\[
\frac{1}{x+1}\in k[x]_{(x)}.
\]

\subsubsection{Exercise 3}

Determine whether
\[
\frac{x^2+1}{x-2}
\]
belongs to
\[
k[x]_{(x)}
\qquad\text{and}\qquad
k[x]_{(x-2)}.
\]

\textbf{Solution.}

For
\[
k[x]_{(x)},
\]
the denominator must satisfy
\[
(x-2)(0)=-2\neq 0.
\]
So
\[
x-2\notin (x),
\]
hence
\[
\frac{x^2+1}{x-2}\in k[x]_{(x)}.
\]

For
\[
k[x]_{(x-2)},
\]
the denominator must not lie in \((x-2)\). But
\[
x-2\in (x-2),
\]
so
\[
\frac{x^2+1}{x-2}\notin k[x]_{(x-2)}.
\]

\subsubsection{Exercise 4}

Determine whether
\[
\frac{x^2+1}{x^2-1}
\]
belongs to
\[
k[x]_{(x-1)}.
\]

\textbf{Solution.}

The denominator is
\[
x^2-1=(x-1)(x+1).
\]
Since it vanishes at \(x=1\),
\[
x^2-1\in (x-1).
\]
Therefore
\[
\frac{x^2+1}{x^2-1}\notin k[x]_{(x-1)}.
\]

\subsubsection{Exercise 5}

Determine whether
\[
\frac{x^2+1}{x^2-1}
\]
belongs to
\[
k[x]_{(x)}.
\]

\textbf{Solution.}

We test the denominator at \(x=0\):
\[
x^2-1\big|_{x=0}=-1\neq 0.
\]
So
\[
x^2-1\notin (x),
\]
and hence
\[
\frac{x^2+1}{x^2-1}\in k[x]_{(x)}.
\]

\subsubsection{Exercise 6}

Determine whether
\[
\frac{x^3+x+1}{x^5-7x+2}
\]
belongs to
\[
k[x]_{(0)}.
\]

\textbf{Solution.}

Since
\[
k[x]_{(0)}=k(x),
\]
every fraction with nonzero polynomial denominator belongs to this localization. So if
\[
x^5-7x+2\neq 0
\]
as a polynomial, then
\[
\frac{x^3+x+1}{x^5-7x+2}\in k[x]_{(0)}.
\]

\subsubsection{Exercise 7}

Determine whether
\[
\frac{1}{x(x+1)}
\]
belongs to
\[
k[x]_x
\qquad\text{and}\qquad
k[x]_{(0)}.
\]

\textbf{Solution.}

In
\[
k[x]_x,
\]
the denominator must be a power of \(x\). But
\[
x(x+1)
\]
is not a power of \(x\), so
\[
\frac{1}{x(x+1)}\notin k[x]_x.
\]

In
\[
k[x]_{(0)}=k(x),
\]
every nonzero polynomial denominator is allowed, so
\[
\frac{1}{x(x+1)}\in k[x]_{(0)}.
\]

\subsection{Exercises in two variables}

\subsubsection{Exercise 8}

Determine whether
\[
\frac{1}{x}
\]
belongs to
\[
k[x,y]_{(x,y)}
\qquad\text{and}\qquad
k[x,y]_{(y)}.
\]

\textbf{Solution.}

In
\[
k[x,y]_{(x,y)},
\]
the denominator must not vanish at \((0,0)\). But
\[
x(0,0)=0,
\]
so
\[
x\in (x,y),
\]
and therefore
\[
\frac{1}{x}\notin k[x,y]_{(x,y)}.
\]

In
\[
k[x,y]_{(y)},
\]
the denominator is allowed if it is not divisible by \(y\). Since
\[
x\notin (y),
\]
we get
\[
\frac{1}{x}\in k[x,y]_{(y)}.
\]

\subsubsection{Exercise 9}

Determine whether
\[
\frac{1}{x+y}
\]
belongs to
\[
k[x,y]_{(x,y)}
\qquad\text{and}\qquad
k[x,y]_{(y)}.
\]

\textbf{Solution.}

At the origin,
\[
(x+y)(0,0)=0,
\]
so
\[
x+y\in (x,y).
\]
Hence
\[
\frac{1}{x+y}\notin k[x,y]_{(x,y)}.
\]

Now check whether
\[
x+y\in (y).
\]
This is false, because \(x+y\) is not divisible by \(y\). Therefore
\[
\frac{1}{x+y}\in k[x,y]_{(y)}.
\]

\subsubsection{Exercise 10}

Determine whether
\[
\frac{x^2+y}{1+x+y}
\]
belongs to
\[
k[x,y]_{(x,y)}.
\]

\textbf{Solution.}

We evaluate the denominator at the origin:
\[
(1+x+y)(0,0)=1.
\]
So
\[
1+x+y\notin (x,y).
\]
Hence
\[
\frac{x^2+y}{1+x+y}\in k[x,y]_{(x,y)}.
\]

\subsubsection{Exercise 11}

Determine whether
\[
\frac{x}{y}
\]
belongs to
\[
k[x,y]_{(y)}
\qquad\text{and}\qquad
k[x,y]_y.
\]

\textbf{Solution.}

In
\[
k[x,y]_{(y)},
\]
we invert all elements not in \((y)\). But
\[
y\in (y),
\]
so
\[
\frac{x}{y}\notin k[x,y]_{(y)}.
\]

In
\[
k[x,y]_y,
\]
we explicitly invert powers of \(y\). Since the denominator is \(y\),
\[
\frac{x}{y}\in k[x,y]_y.
\]

\subsubsection{Exercise 12}

Determine whether
\[
\frac{1}{xy}
\]
belongs to
\[
k[x,y]_x,
\qquad
k[x,y]_y,
\qquad
k[x,y]_{(x,y)}.
\]

\textbf{Solution.}

In
\[
k[x,y]_x,
\]
only powers of \(x\) are allowed as denominators. Since
\[
xy
\]
is not a power of \(x\),
\[
\frac{1}{xy}\notin k[x,y]_x.
\]

Similarly,
\[
\frac{1}{xy}\notin k[x,y]_y,
\]
because \(xy\) is not a power of \(y\).

In
\[
k[x,y]_{(x,y)},
\]
the denominator must not vanish at \((0,0)\). But
\[
xy(0,0)=0,
\]
so
\[
\frac{1}{xy}\notin k[x,y]_{(x,y)}.
\]

\subsubsection{Exercise 13}

Determine whether
\[
\frac{1}{1+xy}
\]
belongs to
\[
k[x,y]_{(x,y)}.
\]

\textbf{Solution.}

We compute
\[
(1+xy)(0,0)=1.
\]
Thus
\[
1+xy\notin (x,y),
\]
so
\[
\frac{1}{1+xy}\in k[x,y]_{(x,y)}.
\]

\subsubsection{Exercise 14}

Determine whether
\[
\frac{1}{y^2+x}
\]
belongs to
\[
k[x,y]_{(y)}.
\]

\textbf{Solution.}

To belong to
\[
k[x,y]_{(y)},
\]
the denominator must not lie in \((y)\). The polynomial
\[
y^2+x
\]
is not divisible by \(y\), since it has the term \(x\). Therefore
\[
y^2+x\notin (y),
\]
and hence
\[
\frac{1}{y^2+x}\in k[x,y]_{(y)}.
\]

\subsection{Exercises comparing quotient and localization}

\subsubsection{Exercise 15}

Is
\[
\frac{1}{x}
\]
an element of
\[
k[x]/(x)?
\]

\textbf{Solution.}

No. The ring
\[
k[x]/(x)
\]
is just
\[
k.
\]
In this quotient one has
\[
x=0,
\]
so the symbol
\[
\frac{1}{x}
\]
does not make sense there. Quotients do not create fractions; they only impose relations. Thus
\[
\frac{1}{x}\notin k[x]/(x).
\]

\subsubsection{Exercise 16}

Compare
\[
\frac{1}{x},
\qquad
\frac{1}{x+1}
\]
in
\[
k[x]_x,
\qquad
k[x]_{(x)},
\qquad
k[x]_{(0)}.
\]

\textbf{Solution.}

For
\[
\frac{1}{x}:
\]
\[
\frac{1}{x}\in k[x]_x,
\qquad
\frac{1}{x}\notin k[x]_{(x)},
\qquad
\frac{1}{x}\in k[x]_{(0)}.
\]

For
\[
\frac{1}{x+1}:
\]
\[
\frac{1}{x+1}\notin k[x]_x,
\qquad
\frac{1}{x+1}\in k[x]_{(x)},
\qquad
\frac{1}{x+1}\in k[x]_{(0)}.
\]

This exercise clearly shows the differences among the three localizations.

\subsection{A general test}

The preceding exercises illustrate the basic criterion.

\subsubsection{Exercise 17}

Let
\[
A=k[x],
\qquad
\mathfrak p=(x-a).
\]
Show that
\[
\frac{f(x)}{g(x)}\in A_{\mathfrak p}
\]
if and only if
\[
g(a)\neq 0.
\]

\textbf{Solution.}

By definition,
\[
A_{\mathfrak p}=S^{-1}A,
\qquad
S=A\setminus \mathfrak p.
\]
So \(\frac{f(x)}{g(x)}\) belongs to \(A_{\mathfrak p}\) exactly when
\[
g(x)\notin \mathfrak p=(x-a).
\]
By the factor theorem,
\[
g(x)\in (x-a)
\quad\Longleftrightarrow\quad
g(a)=0.
\]
Therefore
\[
g(x)\notin (x-a)
\quad\Longleftrightarrow\quad
g(a)\neq 0.
\]
Hence
\[
\frac{f(x)}{g(x)}\in k[x]_{(x-a)}
\quad\Longleftrightarrow\quad
g(a)\neq 0.
\]

\subsubsection{Exercise 18}

Let
\[
A=k[x,y],
\qquad
\mathfrak m=(x-a,y-b).
\]
Show that
\[
\frac{f(x,y)}{g(x,y)}\in A_{\mathfrak m}
\]
if and only if
\[
g(a,b)\neq 0.
\]

\textbf{Solution.}

By definition,
\[
A_{\mathfrak m}=S^{-1}A,
\qquad
S=A\setminus \mathfrak m.
\]
Thus \(\frac{f}{g}\) belongs to \(A_{\mathfrak m}\) exactly when
\[
g\notin \mathfrak m.
\]
But
\[
\mathfrak m=(x-a,y-b)
\]
consists exactly of those polynomials vanishing at \((a,b)\). Therefore
\[
g\notin (x-a,y-b)
\quad\Longleftrightarrow\quad
g(a,b)\neq 0.
\]
So
\[
\frac{f(x,y)}{g(x,y)}\in k[x,y]_{(x-a,y-b)}
\quad\Longleftrightarrow\quad
g(a,b)\neq 0.
\]

\subsection{Final summary}

In polynomial rings, localization is best understood by asking which denominators are allowed.

For
\[
k[x]_f,
\]
the denominator must be a power of \(f\).

For
\[
k[x]_{(x-a)},
\]
the denominator must not vanish at \(x=a\).

For
\[
k[x,y]_{(x-a,y-b)},
\]
the denominator must not vanish at \((a,b)\).

For
\[
k[x]_{(0)},
\]
every nonzero polynomial is allowed, so one gets the field of rational functions
\[
k(x).
\]

The worked exercises above show that most membership questions in localizations reduce to one simple test:

\[
\text{Is the denominator in the multiplicative set, or not?}
\]

\end{document}

